\section{Dataset}
\label{sec:dataset}

The dataset used in this work is organized as a source-retrieval task in which, given a social-media post expressing a
scientific claim and implicitly referring to a scientific publication without linking to it explicitly, the system is
required to identify the corresponding paper within a fixed pool of candidate documents. The benchmark is multilingual
on the query side, as posts are provided in English, German, and French, whereas all queries are matched against the
same shared collection of scientific papers.

The document collection consists of 10,000 scientific publications, each identified by a unique \texttt{pubkey} and
described through the fields \texttt{title}, \texttt{abstract}, \texttt{venue}, and \texttt{authors}; in our pipeline,
each candidate document is represented by the concatenation of title and abstract, a choice that provides a compact yet
informative textual representation for both retrieval and reranking, while preserving a uniform document space across
all languages.

The query data are organized by language and by split: for English, German, and French, both the training and
development portions contain the query identifier (\texttt{index}), the post text, and the gold \texttt{pubkey} of the
referenced paper, thus supporting supervised training as well as model selection and internal evaluation, the released
test files, instead, include only the query identifier and the post text, without the corresponding gold document
identifier.

The Table~\ref{tab:dataset-stats} reports the size of all the data splits considered in our experiments.

\begin{center}
    \captionof{table}{Dataset statistics for the version used in our experiments.}
    \label{tab:dataset-stats}
    \begin{tabular}{lcccc}
        \toprule
        Split & EN & DE & FR & Total \\
        \midrule
        Train & 14,977 & 1,460 & 2,807 & 19,244 \\
        Dev   & 3,905  & 386   & 702   & 4,993 \\
        Test  & 6,076  & 876   & 1,220 & 8,172 \\
        \bottomrule
    \end{tabular}
\end{center}

This setting is particularly challenging because the queries take the form of short, noisy, and informal social-media
posts, while the target documents are substantially longer scientific texts written in a more formal and
information-dense style. As a result, the system must bridge a marked gap in both length and linguistic register between
the two sides of the retrieval problem; moreover, the multilingual nature of the benchmark further increases the
difficulty, since queries in different languages must all be matched against the same shared document collection,
requiring the model to combine semantic matching with cross-lingual generalization.
